\documentclass[12pt,a4paper]{article}

\usepackage[utf8]{inputenc}
\usepackage[russian]{babel}
\usepackage[margin=2.5cm]{geometry}
\usepackage{amsmath,amssymb}
\usepackage{booktabs}
\usepackage{parskip}
\usepackage[hidelinks]{hyperref}

\title{\textbf{Физически обоснованное раннее предупреждение\\
о квантовой декогеренции при неизвестной\\
нестационарной скорости диссипации}\\[0.5em]
}

\date{}

\begin{document}
\maketitle

\begin{center}
\fbox{\parbox{0.9\textwidth}{\small
\textbf{Статус документа.} Это пояснительный черновик на русском языке, а не
публикация. Авторитетные источники результатов — статьи
\texttt{paper/paper\_1/main.tex} (1 кубит, сценарии A--C) и
\texttt{paper/paper\_2/paper2.tex} (2 кубита, сценарии D/E); при расхождении
верны они. Числа ниже сверены с
\texttt{experiments/results/*.csv} по состоянию на сентябрь 2026; у каждой
таблицы указан эксперимент-источник.
}}
\end{center}

\tableofcontents
\newpage

% ============================================================
\section{Проблема}

Квантовые компьютеры, квантовые сенсоры и системы квантовой связи
принципиально ограничены явлением \textbf{декогеренции} —
необратимой потерей квантовой когерентности из-за взаимодействия
системы с окружающей средой.

Математически открытая квантовая система описывается уравнением
Линдблада:
\begin{equation}
  \dot{\rho} = -i[H,\rho]
  + \gamma(t)\!\sum_k\!\left(
    L_k\rho L_k^\dagger
    - \tfrac{1}{2}\{L_k^\dagger L_k,\rho\}
  \right),
\end{equation}
где $\rho$ — матрица плотности, $H$ — гамильтониан, $L_k$ — операторы
прыжков, а $\gamma(t) \geq 0$ — нестационарная скорость диссипации.

\textbf{Оператор прыжка $\sigma_-$.}
Для двухуровневой системы $\{|0\rangle, |1\rangle\}$ оператор
амплитудного затухания (канал релаксации) задаётся матрицей:
\[
  \sigma_- = |0\rangle\langle 1|
  = \begin{pmatrix} 0 & 1 \\ 0 & 0 \end{pmatrix}.
\]
Он переводит кубит из возбуждённого состояния $|1\rangle$ в основное
$|0\rangle$ — моделирует спонтанное излучение или релаксацию в
термостат при нулевой температуре.
Подставляя $L = \sigma_-$ в уравнение Линдблада и раскрывая скобки:
\[
  \sigma_-\rho\sigma_+
  - \tfrac{1}{2}\bigl(\sigma_+\sigma_-\rho + \rho\,\sigma_+\sigma_-\bigr),
  \quad \sigma_+ = (\sigma_-)^\dagger = \begin{pmatrix} 0 & 0 \\ 1 & 0 \end{pmatrix},
\]
получаем уравнения на элементы $\rho$:
\[
  \dot\rho_{11} = -\gamma\rho_{11},
  \qquad
  \dot\rho_{01} = -\tfrac{\gamma}{2}\,\rho_{01}.
\]
Отсюда $\rho_{01}(t) = \rho_{01}(0)\,e^{-\gamma t/2}$, и когерентность
убывает как $C(t) \propto e^{-\gamma t/2}$.

Для фазовой релаксации используется другой оператор — $\sigma_z$
(диагональная матрица $\mathrm{diag}(1,-1)$), который не меняет
населённости, но уничтожает когерентность вдвое быстрее:
$\dot\rho_{01} = -\gamma\,\rho_{01}$, то есть $T_2 = 1/\gamma$.

\textbf{Время декогеренции $T_2$} — это момент, когда
$\ell_1$-когерентность (сумма модулей внедиагональных элементов
матрицы плотности) падает до $e^{-1}$ от начального значения.

Знание $T_2$ наперёд — или, точнее, оценка \emph{оставшегося}
времени до декогеренции из неполной ранней траектории — напрямую
полезно: оно позволяет системам квантовой коррекции ошибок заблаговременно
сбросить или перекодировать кубит до критического падения верности.

% ============================================================
\section{Когда задача нетривиальна}

\subsection{Случай постоянного \texorpdfstring{$\gamma$}{гамма}: аналитическое решение}

Здесь $C(t)$ — $\ell_1$-когерентность матрицы плотности
(сумма модулей всех внедиагональных элементов $\sum_{i\neq j}|\rho_{ij}|$);
для однокубитной системы это просто $2|\rho_{01}(t)|$.
При $C = 0$ система полностью декогерирована.

Если $\gamma$ постоянна и используется канал амплитудного затухания
($\sigma_-$), когерентность убывает экспоненциально:
\[
  C(t) = C_0\,e^{-\gamma t/2}
  \;\Longrightarrow\;
  \frac{d\log C}{dt} = -\frac{\gamma}{2}
  \;\Longrightarrow\;
  T_2 = \frac{2}{\gamma} = -\frac{1}{\mathrm{slope}},
\]
где slope — наклон $\log C(t)$ по МНК на любом участке траектории до
декогеренции.

Мы проверили: эта формула даёт $R^2 = 0{,}9995$ и
MAE~$= 0{,}056$ на тестовой выборке.
\textbf{Машинное обучение здесь не нужно.}

\subsection{Случай переменного \texorpdfstring{$\gamma(t)$}{гамма(t)}: реалистичная ситуация}

На любой физической платформе скорость диссипации нестационарна:

\begin{center}
\begin{tabular}{p{3.8cm}p{7cm}}
\toprule
Платформа & Источник нестационарного $\gamma(t)$ \\
\midrule
Сверхпроводящие кубиты &
  $1/f$-флуктуации потока, TLS-флуктуаторы \\
Спиновые кубиты &
  Сверхтонкое взаимодействие с ядерным спин-резервуаром \\
Ионные ловушки &
  Дрейф мощности лазера, флуктуации микродвижения \\
Фотонные системы &
  Дрейф скорости утечки из резонатора \\
\bottomrule
\end{tabular}
\end{center}

В этих условиях наклон МНК даёт лишь \emph{локальную} оценку текущего
$\gamma(t)/2$, и предсказанное $T_2$ может быть систематически
смещено — в зависимости от формы профиля $\gamma(t)$
(монотонный рост/спад, пик, ступенька и т.\,д.).

Задача становится принципиально нетривиальной:
никакой аналитической формулы общего вида нет.
Именно здесь нужен обученный предиктор.

% ============================================================
\section{Отличие от существующих подходов}

Ближайшая родственная работа — arXiv:2505.06928
(\textit{Phys.~Rev.~A}, 2025) — использует Transformer для
\textbf{обратной задачи}: восстановления $\gamma(t)$ из
\textit{полной} траектории.
Это офлайн-подход, неприменимый в режиме реального времени.

\begin{center}
\begin{tabular}{p{3cm}p{3.5cm}p{4cm}}
\toprule
Характеристика & arXiv:2505.06928 & \textbf{Наша работа} \\
\midrule
Цель & Восстановить $\gamma(t)$ & Предсказать $T_2$ \\
Тип задачи & Обратная (офлайн) & Раннее предупреждение (онлайн) \\
Вход & Полная траектория & Ранний частичный отрезок \\
Выход & Коэффициенты Бернштейна & $T_2 +$ score риска \\
Физика & Отсутствует & МНК-наклон встроен как prior \\
Аналитический предел & Нет & Восстанавливает формулу при $\gamma$=const \\
\bottomrule
\end{tabular}
\end{center}

Наш подход — \textbf{прямой и онлайновый}: предсказание из первых
$\approx 15$--$20\%$ траектории, без необходимости знать $\gamma(t)$
или запускать полное моделирование.

% ============================================================
\section{Архитектура: Physics-Informed Residual BiLSTM}

\subsection{Ключевая идея}

Сеть не предсказывает $T_2$ с нуля.
Она предсказывает \textbf{поправку} $\delta T$ к аналитической оценке:
\begin{equation}
  \widehat{T}_{2,\text{rem}}
  = \underbrace{T_{\text{physics}}(\mathbf{w})}_{\text{МНК-оценка}}
  +\; \underbrace{\delta T_{\text{LSTM}}(\mathbf{w})}_{\text{выученная коррекция}}.
\end{equation}

Это принципиально иное, чем классические PINN (Raissi et al., 2019),
где физика входит как \emph{штраф в функцию потерь} и теоретически
может быть проигнорирована достаточно гибкой сетью.
Здесь физика входит как \textbf{архитектурный inductive bias}:
выход сети структурно декомпозирован на физическую часть и
выученную поправку.
Это гарантирует:
\begin{enumerate}
  \item \textbf{Graceful degradation}: при $\gamma = \mathrm{const}$
    физическая формула точна, поэтому сеть учит $\delta T \approx 0$
    и вся система деградирует к аналитическому решению.
  \item \textbf{Corrective capacity}: при переменном $\gamma(t)$
    физический prior даёт смещённую, но сильную оценку —
    сеть корректирует только остаток.
    Это намного проще, чем предсказывать $T_2$ напрямую.
\end{enumerate}

\subsection{Детали архитектуры}

\textbf{Вход.} Временной ряд наблюдаемых за окно длины $W = 20$ шагов:
$\langle\sigma_x\rangle$, $\langle\sigma_y\rangle$,
$\langle\sigma_z\rangle$, $\ell_1$-когерентность, чистота — плюс
их логарифмические преобразования ($F = 7$ для 1-кубитных,
$F = 10$ для 2-кубитных систем после нулевого дополнения).

\textbf{Энкодер.} Двунаправленный LSTM (hidden~64, 1 слой) →
mean-pooling → вектор размерности 128.

\textbf{Физические скаляры.} Четыре числа вычисляются аналитически
по окну: наклон МНК $(\approx -\gamma/2)$, $\log C$ в конце окна,
текущее время $t_\mathrm{obs}$, оценка $T_\mathrm{physics}$.
Конкатенируются к LSTM-эмбеддингу: итого 132 измерения.

\textbf{Головы.}
\begin{itemize}
  \item \textit{Regression head}: Linear(64) $\to$ GELU $\to$ Linear(1)
    $\to \delta T$. Итог: $T_\mathrm{physics} + \delta T$.
  \item \textit{Risk head}: Linear(32) $\to$ GELU $\to$ Linear(1)
    $\to$ sigmoid $\to p_\mathrm{risk}$.
\end{itemize}

\textbf{Score риска $p_\mathrm{risk}$.}
Помимо регрессии $T_2$, сеть параллельно решает задачу бинарной
классификации: грозит ли декогеренция в ближайшем будущем?
Метка $y_\mathrm{risk} \in \{0, 1\}$ формируется автоматически из симуляции:
$y_\mathrm{risk} = 1$, если оставшееся время до декогеренции меньше порогового.
Выход $p_\mathrm{risk} \in [0,1]$ — вероятность этого события,
получаемая через сигмоиду из выхода \textit{risk head}.

\textbf{Обучение.} Совместная функция потерь:
\begin{equation}
  \mathcal{L}
  = \underbrace{\mathcal{L}_{\mathrm{Huber}}(\delta T,\,\delta T^*)}_{\text{регрессия поправки}}
  + \underbrace{\mathcal{L}_{\mathrm{BCE}}(p_\mathrm{risk},\,y_\mathrm{risk})}_{\text{бинарная классификация}},
\end{equation}
где $\delta T^* = T_2^* - T_{\mathrm{physics}}$ — истинная поправка (цель регрессии).

\textbf{Huber loss} (робастная регрессия):
\[
  \mathcal{L}_{\mathrm{Huber}}(x, x^*) =
  \begin{cases}
    \tfrac{1}{2}(x - x^*)^2, & |x - x^*| \leq \delta, \\
    \delta\!\left(|x - x^*| - \tfrac{\delta}{2}\right), & |x - x^*| > \delta.
  \end{cases}
\]
При малых ошибках ведёт себя как MSE (квадратичный штраф),
при больших — как MAE (линейный), что снижает влияние выбросов
по сравнению с чистым MSE.

\textbf{BCE} (бинарная кросс-энтропия):
\[
  \mathcal{L}_{\mathrm{BCE}}(p, y)
  = -\,y\log p - (1 - y)\log(1 - p).
\]
Стандартная функция потерь для задач бинарной классификации:
штрафует за уверенные неправильные предсказания сильнее, чем за
неуверенные.

Оптимизатор AdamW, планировщик reduce-on-plateau,
ранняя остановка по валидационному лоссу (patience = 25 эпох).

\subsection{Ablation-варианты (эксперимент E1)}

\begin{center}
\begin{tabular}{ll}
\toprule
Вариант & Описание \\
\midrule
\texttt{physics\_only} &
  Только формула МНК, без нейронной сети \\
\texttt{lstm\_only} &
  BiLSTM без физического prior ($T_\mathrm{physics} = 0$) \\
\texttt{physics\_lstm} &
  Полная residual-архитектура \textbf{(наш метод)} \\
\texttt{transformer} &
  Чистый Transformer, аналог arXiv:2505.06928 \\
\bottomrule
\end{tabular}
\end{center}

% ============================================================
\section{Физические системы и параметризация \texorpdfstring{$\gamma(t)$}{гамма(t)}}

\subsection{Поддерживаемые системы}

\begin{center}
\begin{tabular}{p{1.2cm}p{2cm}p{4cm}p{3cm}}
\toprule
Сцен. & Система & Гамильтониан & Операторы прыжков \\
\midrule
A & 1q & $H = \omega\sigma_z$ &
  $\sigma_-$, $\gamma{=}\mathrm{const}$ \\
B & 1q & $H = \omega\sigma_z$ &
  $\sigma_-$, TD $\gamma(t)$ \\
C & 1q & $H = \omega\sigma_z$ &
  $\sigma_z$, TD $\gamma(t)$ \\
D & 2q XXZ & $J(\sigma^x\sigma^x{+}\sigma^y\sigma^y{+}\sigma^z\sigma^z)$ &
  $\sigma_-$ на каждый \\
E & 2q TFIM & $J\sigma^x\sigma^x{+}h(\sigma^z{+}\sigma^z)$ &
  $\sigma_-$ на каждый \\
\bottomrule
\end{tabular}
\end{center}

\subsection{Параметризация \texorpdfstring{$\gamma(t)$}{гамма(t)} полиномами Бернштейна}

Для гарантированно неотрицательного $\gamma(t) \geq 0$ при любых
параметрах используются полиномы Бернштейна 4-й степени с
неотрицательными коэффициентами.
В обучающей выборке используются шесть качественно различных форм:
постоянная, монотонно возрастающая/убывающая, пик, ступенька вверх/вниз.

% ============================================================
\section{Эксперименты и результаты}

\subsection{E1 — Ablation: компоненты архитектуры}

Одно- и двухкубитные сценарии измерялись \emph{разными} экспериментами с
разными настройками, поэтому приводятся отдельно — объединять их в одну
таблицу было бы некорректно.

\textbf{(а) Сценарии A--C} — эксперимент E1 (100 траекторий на сценарий,
источник: \texttt{e1\_ablation\_p1.csv}; Paper 1, табл.~E1):

\begin{center}
\begin{tabular}{lrrr}
\toprule
Вариант & A & B & C \\
\midrule
\texttt{physics\_only}  & \textbf{0{,}999} & 0{,}753 & 0{,}965 \\
\texttt{lstm\_only}     & $-$0{,}446 & 0{,}009 & $-$0{,}029 \\
\texttt{physics\_lstm}  & 0{,}976 & \textbf{0{,}945} & 0{,}976 \\
\texttt{transformer}    & 0{,}799 & 0{,}902 & \textbf{0{,}981} \\
\bottomrule
\end{tabular}
\end{center}

\textbf{(б) Сценарии D, E} — эксперимент E6a (500 траекторий на сценарий,
adaptive prior $\tau = 0{,}7$; источник:
\texttt{e6\_ablation\_improved.csv}; Paper 2, табл.~E6a):

\begin{center}
\begin{tabular}{lrr}
\toprule
Вариант & D (XXZ) & E (TFIM) \\
\midrule
\texttt{physics\_only}  & $-$0{,}687 & $-$4{,}267 \\
\texttt{lstm\_only}     & 0{,}673 & 0{,}377 \\
\texttt{physics\_lstm}  & 0{,}636 & 0{,}455 \\
\texttt{transformer}    & \textbf{0{,}897} & \textbf{0{,}689} \\
\bottomrule
\end{tabular}
\end{center}
(Метрика: $R^2$, выше — лучше.)

\textbf{Выводы:}
\begin{itemize}
  \item \textbf{Сценарий A}: \texttt{physics\_only} $R^2 = 0{,}999$ —
    аналитическая формула практически точна при постоянном $\gamma$.
    ML не добавляет ничего. Это \emph{подтверждает} физическую
    состоятельность prior.
  \item \textbf{Сценарии B, C}: \texttt{physics\_lstm} превосходит
    \texttt{physics\_only} на B ($0{,}945$ vs $0{,}753$) и слегка — на C
    ($0{,}976$ vs $0{,}965$). Ключевое: \texttt{lstm\_only} проваливается
    на всех трёх сценариях ($-0{,}446$, $0{,}009$, $-0{,}029$), то есть
    физический prior \emph{необходим}, а не просто полезен.
  \item \textbf{Сценарии D, E}: OLS-формула даёт $R^2 < 0$
    ($-0{,}687$ и $-4{,}267$) — для 2-кубитных систем когерентность
    не является простой экспонентой из-за запутанности.
    Adaptive prior выправляет ситуацию:
    \texttt{physics\_lstm} даёт $0{,}636$ (D) и $0{,}455$ (E).
  \item \textbf{Transformer доминирует на двух кубитах} ($0{,}897$ на D
    и $0{,}689$ на E) — в отличие от однокубитного случая, где
    \texttt{physics\_lstm} сильнее. Там, где физический prior
    структурно неверен, гибкость self-attention выгоднее; это
    зафиксировано в Paper 2 как \emph{опровержение} переноса
    residual-архитектуры на 2 кубита.
\end{itemize}

\subsection{E2 — Минимальное окно наблюдения}

\begin{center}
\begin{tabular}{rrrr}
\toprule
$f = t_\mathrm{obs}/T_2$ & $R^2$ & MAE & AUROC \\
\midrule
0{,}05 & 0{,}711 & 0{,}825 & 0{,}933 \\
0{,}10 & 0{,}723 & 0{,}813 & 0{,}950 \\
0{,}15 & 0{,}735 & 0{,}773 & 0{,}947 \\
0{,}20 & 0{,}736 & 0{,}702 & 0{,}872 \\
0{,}30 & 0{,}734 & 0{,}677 & 0{,}969 \\
0{,}50 & 0{,}697 & 0{,}674 & 0{,}990 \\
\bottomrule
\end{tabular}
\end{center}
(Источник: \texttt{e2\_window\_sweep\_p1.csv}; смесь сценариев B+C.)

Примечательно: \textbf{AUROC~$\geq 0{,}93$ уже при $f = 0{,}05$} —
сигнал тревоги о скорой декогеренции работает уже через 5\% траектории.
$R^2$ остаётся примерно постоянным ($0{,}70$--$0{,}74$) при всех $f$:
узкое место — ёмкость обучения и смешение режимов B и C, а не
информационное содержание окна (на одном сценарии B модель достигает
$0{,}945$ при том же $W = 20$).
Заявлять $R^2 \geq 0{,}95$ по этим данным нельзя — раннее предупреждение
здесь является \emph{классификационным} результатом (AUROC), а не
регрессионным.

\subsection{E3 — Устойчивость к шуму измерений}

\begin{center}
\begin{tabular}{rrrr}
\toprule
$\sigma$ & $R^2$ & MAE & AUROC \\
\midrule
0{,}000 & 0{,}959 & 0{,}220 & 0{,}976 \\
0{,}010 & 0{,}956 & 0{,}224 & 0{,}976 \\
0{,}020 & 0{,}949 & 0{,}251 & 0{,}971 \\
0{,}050 & 0{,}901 & 0{,}374 & 0{,}942 \\
0{,}100 & 0{,}727 & 0{,}597 & 0{,}865 \\
\bottomrule
\end{tabular}
\end{center}
(Источник: \texttt{e3\_noise\_sweep\_p1.csv}, $N = 74$.)

При реалистичном уровне шума $\sigma = 0{,}05$ модель держит
$R^2 = 0{,}901$ и AUROC $= 0{,}942$.
Деградация от $\sigma = 0$ до $\sigma = 0{,}10$:
$\Delta R^2 = -0{,}232$, $\Delta\mathrm{AUROC} = -0{,}111$ — при
$\sigma = 0{,}10$ AUROC опускается до $0{,}865$, что мы считаем
практическим потолком шумоустойчивости текущей архитектуры.
\textbf{Физический prior обеспечивает встроенную устойчивость к шуму}:
OLS-наклон усредняет шум по $W$ точкам ($\sim\!\sqrt{W}$ подавление),
BiLSTM дополнительно использует temporal pooling.

\subsection{E4 — Кросс-системное обобщение}

Единой модели сразу на все пять сценариев в текущих результатах нет.
Кросс-системное обобщение меряется двумя отдельными прогонами:
на однокубитной смеси (A, B, C) и на двухкубитной (D, E).

\textbf{(а) Смесь A+B+C}, \texttt{physics\_lstm}
(источник: \texttt{e4\_cross\_system\_p1.csv}; Paper 1, табл.~E4):

\begin{center}
\begin{tabular}{llrrr}
\toprule
Сценарий & Система & $R^2$ & MAE & AUROC \\
\midrule
A & 1q $\sigma_-$, $\gamma$=const & 0{,}490 & 1{,}164 & 0{,}979 \\
B & 1q $\sigma_-$, TD $\gamma$    & 0{,}867 & 0{,}397 & 0{,}981 \\
C & 1q $\sigma_z$, TD $\gamma$    & \textbf{0{,}936} & \textbf{0{,}162} & \textbf{0{,}990} \\
\bottomrule
\end{tabular}
\end{center}

\textbf{(б) Смесь D+E}, Transformer, эксперимент E11a
(окно $L = 20$, min-gap $= 10$, 500 траекторий на сценарий; Paper 2):

\begin{center}
\begin{tabular}{llrrr}
\toprule
Сценарий & Система & $R^2$ & MAE & AUROC \\
\midrule
D & 2q XXZ, TD $\gamma$  & \textbf{0{,}817} & 0{,}616 & \textbf{0{,}937} \\
E & 2q TFIM, TD $\gamma$ & 0{,}566 & 0{,}900 & 0{,}888 \\
\bottomrule
\end{tabular}
\end{center}

Одна модель покрывает свою группу систем без переобучения.
На сценарии A модель хуже специализированного аналитического метода
($0{,}490$ vs $0{,}999$) — ожидаемо при обобщении: при постоянном $\gamma$
физический член уже точен, и обученная на смеси сеть зря добавляет
ненулевую поправку $\delta T$.
На TFIM (E) $R^2 = 0{,}566$ — цель $0{,}70$ не взята. В Paper 2 это
объясняется структурным потолком, а не недообучением: диссипатор не коммутирует
с оператором чётности и потому мешает два симметрийных сектора, тем сильнее,
чем ближе их временные масштабы. Формулировку «потолок \emph{вблизи фазового
перехода}» пришлось уточнить: репликация J-скана на трёх сидах (R5) не
подтвердила провал, локализованный у $J\approx h$, — минимум кочует по сидам, а
сама точка перехода ничем не выделяется. Устоял монотонный спад
предсказуемости с ростом связи, чего теория и требует: отношение расщеплений
$r(J)$ растёт монотонно и пика в $J=h$ не имеет.

% ============================================================
\section{Место в литературе и научный вклад}

\subsection{Отличие от PINNs}

Классические PINN \cite{raissi2019} включают физику в функцию потерь
в виде штрафа за нарушение дифференциального уравнения.
Наш подход принципиально отличается: физика входит как
\textbf{жёсткий архитектурный prior} — выход сети структурно задан
как $T_\mathrm{physics} + \delta T_\mathrm{LSTM}$.
Преимущества:
\begin{itemize}
  \item Физический prior не может быть «проигнорирован» сетью при
    оптимизации.
  \item В аналитически точном пределе ($\gamma = \mathrm{const}$)
    модель гарантированно деградирует к точному решению при
    $\delta T \to 0$.
  \item Сеть работает в значительно сниженном пространстве поиска
    (только остаток), что ускоряет сходимость и снижает потребность
    в данных.
\end{itemize}

\subsection{Вклады работы}

\begin{enumerate}
  \item \textbf{Чёткая граница применимости ML}: показано, что ML
    нужен только при нестационарном или многокубитном режиме;
    при постоянном $\gamma$ + $\sigma_-$ аналитика точна ($R^2 \approx 1$).
  \item \textbf{Residual-архитектура с физическим prior}: встроенная
    МНК-оценка гарантирует аналитический предел и улучшает
    выборочную эффективность.
  \item \textbf{Онлайн-применимость}: эффективный alarm уже при
    $f = 0{,}05$ (AUROC~$\geq 0{,}92$).
  \item \textbf{Встроенная устойчивость к шуму}: OLS-усреднение по
    окну подавляет измерительный шум архитектурно.
  \item \textbf{Калиброванный score риска}: $p_\mathrm{risk}$ напрямую
    пригоден для интеграции в системы квантовой коррекции ошибок.
\end{enumerate}

Ключевые зависимости: \texttt{qutip}, \texttt{torch}, \texttt{numpy},
\texttt{scikit-learn}, \texttt{gradio} (интерактивный UI).

\begin{thebibliography}{9}
\bibitem{raissi2019}
M.~Raissi, P.~Perdikaris, G.~E.~Karniadakis.
Physics-informed neural networks.
\textit{J.~Comput.~Phys.}, 378:686--707, 2019.
\end{thebibliography}

\end{document}
